\documentclass{article}
\usepackage{import}
\import{../../template/}{article-template}
\graphicspath{{../../images/}}
\addbibresource{../../template/library.bib}

\title{Emacs Lisp}
\author{PENG}
\date{Created: 20260815, Version: 20260831}

\begin{document}
\maketitle
\tableofcontents

\section{Introduction}
\subsection{Lisp History}
Lisp (List Processing Language) was first developed in the late 1950s.

Common Lisp (cl-lib), Scheme (Guile)

\section{Lisp Data Types}
\begin{definition}[object]
    A Lisp object is a piece of data.
\end{definition}

\begin{definition}[primitive types]
    integer, float, cons, symbol, string, vector, hash-table, subr, byte-code function, record, buffer
\end{definition}

\begin{remark}
    The primitive type of each object is \textbf{implicit} in the object itself.
\end{remark}

\subsection{Printed Representation and Read Syntax}
\begin{definition}[printed representation and read syntax]
    \verb|prin1| and \verb|read|
\end{definition}

\begin{remark}
    In most cases, an object's printed representation is also a read syntax for the object.
\end{remark}

\begin{example}[objects have no read syntax]
    \verb|#<type varname>| cannot be read.
    \begin{verbatim}
(current-buffer) => #<buffer elisp.tex>
\end{verbatim}
\end{example}

An expression is primarily a Lisp object.

Interactive evaluation process: read the text $\Rightarrow$ produce a Lisp object $\Rightarrow$ evaluate that object. Evaluation and reading are separate activities.

\subsection{Special Read Syntax}
\begin{definition}[special hash notations]~
    \begin{itemize}
        \item \verb|#<>| : objects have no read syntax
        \item \verb|#'| : shortcut for function
    \end{itemize}
\end{definition}

\subsection{Comments}
\begin{definition}[comment]
    \verb|;| continues to the end of line.
\end{definition}

The Lisp reader discards comments.

\subsection{Lisp Programming Types}
Two categories of types:
\begin{itemize}
    \item Lisp programming types:
          \begin{itemize}
              \item integer
              \item floating-point
              \item character
              \item symbol
              \item sequence
              \item array
              \item string
              \item vector
              \item char-table
              \item bool-vector
              \item hash table
              \item function
              \item macro
              \item primitive function
              \item closure
              \item record
              \item type descriptors
              \item type specifiers
              \item autoload
              \item finalizer
          \end{itemize}
    \item editing (unique to Emacs Lisp)
\end{itemize}

\subsubsection{Integer Type}
\begin{definition}[integer]~
    \begin{itemize}
        \item fixnums : small integers, minimum range: 30 bits
        \item bignums : large integers, arbitrary precision
    \end{itemize}
\end{definition}

\begin{example}[compare]~
    \begin{verbatim}
(eql 1 1.) ; or
(= 1 1.)   ; for all numbers
(eq 1 1.)  ; for fixnums
\end{verbatim}
\end{example}

Read syntax for integer: \verb|1. => 1|

\subsubsection{Floating-Point Type}

\section{Variables}
\subsection{Scoping Rules for Variable Bindings}
\subsubsection{Lexical Binding}
\begin{definition}[lexical binding]
    Any reference to the variable must be located textually within the binding construct.
\end{definition}

Lisp evaluator get the value of a variable process:
look first in the lexical environment $\Rightarrow$ if not then look in the dynamic value cell.

\begin{definition}[closure]
    A closure is created when you define a named or anonymous function with lexical binding enabled.
\end{definition}

\begin{example}~
    \begin{verbatim}
(let ((x 0))
  (defun var ()
    (setq x (1+ x))))

(var) => 0
(var) => 1
(var) => 2
\end{verbatim}
\end{example}

\section{Regular Expression}

\section{Automatic Indentation of code}
Automatic indentation \cite{LEWIS_1998_GNU_EMACS_LISP_REFERENCE_MANUAL}:
\begin{itemize}
    \item What is the right indentation of a line?
    \item When to reindent a line?
\end{itemize}

Indentation function:
\begin{itemize}
    \item parsing forward from some safe starting point until the position of interest
    \item parsing backward from the position of interest
\end{itemize}

Generic indentation engine:
\begin{itemize}
    \item CC-mode indentation code
    \item SMIE
    \item Tree-sitter
\end{itemize}

\subsection{Simple Minded Indentation Engine}
\begin{itemize}
    \item operator precedence grammar
    \item unable to detect syntax error
    \item most programming languages cannot be parsed correctly by SMIE
\end{itemize}

\begin{definition}[smie-setup]~
    \label{def:smie-setup}
    \begin{verbatim}
(smie-setup grammar rules-function :forward-token :backward-token)
\end{verbatim}
\end{definition}

\subsection{Tree-sitter}
Parser-based indentation

\section{Language Server Protocol}

\section{Dynamic Module}
% https://www.phimulambda.org/blog/emacs-dynamic-module.html

\section{Canvas}
% https://monadicsheep.org/blog/an-introduction-to-canvas-in-emacs.html

\printbibliography
\end{document}
